\section{Source audit, corrections, and verification status}
\label{sec:audit}

\subsection{The source anchor and the inspection actually performed}
The repository was anchored at commit
\begin{center}\pin,\end{center}
whose timestamp is 11 October 2026 at 18:22:01 UTC. Both directories in
the research request were inspected at that pin. The source record in
\texttt{provenance/sources.json} lists the relevant paths and the level of
inspection, rather than implying a full audit of the repository.

The targeted reading included the canonical main file and preface, the
research-status chapter, the opening 280 lines of the integration chapter,
and the adjacent-tail and research-question sections of the differential
polygamma continuation. The determinantal harmonic-resonance report was
read for its defining ensemble, Gauss product, first-resonance context,
stated simple-pole harmonic-rank result, and further research questions.
Its requested repeated-pole classification and rational-grid extension
are the direct targets of Section~\ref{sec:classification}.

The incoming directory contained six archives at the anchor. The member
list and selected status/readme material of
\emph{Polygamma Reconstruction and Resonance} and the readme of
\emph{Harmonic Tails, Shifts and Twists} were inspected. A combined retrieval
also exposed readme material from \emph{Harmonic Tail Gamma Continuation},
but it was not treated as a full interior audit. The Beta Crossing,
Quadratic Spectrum, and Same-Direction Complex Orders packages were not
read in full in this continuation. No error or absence of a theorem is
attributed to their unread contents. This targeted comparison is enough
to identify the specific source question being answered, but not to prove
an exhaustive novelty claim.

\subsection{Proposed status updates}
The source question about rank-one harmonic pencils should now be split
into a proved classification and remaining arithmetic-depth questions.
Theorem~\ref{thm:classification} supplies the full residue test in the
specified rational-weight model. Theorem~\ref{thm:critical} gives the
particularly short critical-decay characterization, allowing repeated
poles and multiple rational residue classes. The source's simple-pole
construction is thereby extended, not invalidated.

The source question about harmonic rank controlling depth receives the
constructive upper bound in Theorem~\ref{thm:depth}. It should not be
rewritten as a theorem about \emph{minimal} depth. The possible strict
inequality between harmonic rank, determinant degree, and arithmetic depth
remains an important separate issue.

The next-two-particle-resonance question in the older report should be
cross-linked to the current incoming reconstruction/resonance package,
subject to the repository's usual integration review. The present
finite-difference deformation should be listed separately: its simple
pole and its coefficient $c_d(s)$ are not the spectral excess and pole
structure of every higher Gauss resonance.

\subsection{Corrections and safeguards established here}
No newly false analytic theorem was found in the portions of the existing
material inspected for this article. The following are precise
integration safeguards, with explicit examples rather than allegations
about unread sources.

First, the critical-decay characterization requires a \emph{nonzero}
coefficient $\kappa$ in \eqref{eq:critical-decay}. The faster-decaying
weight \eqref{eq:counterexample} has a rank-one geometric pencil and is
not a $d$th finite difference. Replacing exact critical decay by merely
a big-$O$ condition would make the new theorem false.

Second, the escape amplitude is $A_r=c_dV_r$. A counterterm using $V_r$
for all $s$ gives the wrong derivative rows. In particular, the pole-times-
next-coefficient term $b_{k+1}/(k+1)$ in
\eqref{eq:Tjets-closed} must be retained.

Third, the hard-cutoff constant and the spectral Laurent constant differ
by $\partial_sA_r(1,a)$. Equation~\eqref{eq:cutoff-example} gives a
nonzero elementary discrepancy. The earlier source already recognizes
this distinction; our formulas make it explicit for this family.

Fourth, meromorphic continuation does not commute with an isolated
pointwise specialization of a nonnormally convergent family. At $s=0$
the weights vanish, but the continued value is $-1/4$ for two particles.
This also limits the interpretation of identities at nonpositive orders.

Finally, finite harmonic rank and polylogarithmic representation depth
are not arithmetic independence. An explicit $\Li_{2,2}$ term is not
proof that depth two is necessary. Conversely, the simple-pole report's
special ordinary-zeta endpoint closure should not be extrapolated to
all repeated-pole cases without a separate reduction argument.

\subsection{What the supplied verification proves}
The exact suite \texttt{verification/verify\_exact.py} performs
\textbf{1,037 assertions} in rational or symbolic arithmetic. The
machine-readable record is \texttt{results/exact\_checks.json}.
It includes 576 finite endpoint moment identities, 18 direct ordered-sum
Cauchy--Binet comparisons, 18 residue-product checks, 12 single-zeta
matrix identities, independent-order/shift polarization tests, 144
Gamma--harmonic Bell identities, finite-difference residue inversions,
zero-pencil checks, the geometric counterexample, and 20 deliberately
perturbed weights rejected by the rank-one test. It also checks the
explicit cutoff formula, 44 generating-function coefficients, and the
Stieltjes/Bell coefficient rules.

The independent numerical suite records \textbf{29 diagnostics} in
\texttt{results/numerical\_checks.json}. Nine compare determinants of
Mellin-integral moments with the finite-endpoint/zeta expression, including
noninteger and complex orders and shifts. Two test mixed-order/shift
Mellin identities. Four compare direct finite sums with finite boundary
formulas. Four test hard-cutoff constants through finite formulas at
$N=10^{20}$; six test differentiated Stieltjes brackets at $N=10^{18}$.
Three test the explicit polylogarithmic generator, and one tests the
anchored Gamma primitive.

The standard working precision was 60 decimal digits, with increased
precision for cutoff cancellations. The Mellin-determinant relative
residuals in the recorded run are below $6.1\times10^{-57}$.
The largest relative discrepancy across \emph{all} diagnostics is about
$8.83\times10^{-13}$, occurring in a high derivative of a finite-cutoff
approximation. It would therefore be incorrect to describe the whole
suite as agreement to 60 digits. Each row records its actual discrepancy,
tolerance, and cutoff where applicable. These are floating-point
diagnostics, not certified interval enclosures.

The infinite-family conclusions rest on the proofs in this article.
Finite exact checks test implementation and normalization; they do not
replace the general algebraic or analytic arguments. A successful PDF
build likewise verifies document compilation, not mathematical truth.

\subsection{Integration boundary}
The package is self-contained and includes standalone and modular TeX,
the compiled PDF, verification programs and result records, a build
script, a source-scope record, an explicit claim-status file, and checksum
verification. The suggested repository destination is recorded in
\path{integration/INTEGRATION.md}. No repository files were changed,
no archive was retired, and no commit or pull request was created.
The manuscript's Gaussian $S_6$ and revised $S_8$ proof statuses are
unchanged \cite{repo:discovery}.
